\chapter{Discussion}
\label{ch:discussion}

Chapter~\ref{ch:evaluation} reports what was measured. This chapter interprets
it: what the six verdicts mean taken together, how far each estimate can be
pushed, and what the study cannot support. Section~\ref{sec:discussion_core}
states the central finding and Section~\ref{sec:discussion_h1h6} resolves the
tension the pre-registration deliberately built between H1 and H6.
Sections~\ref{sec:external_interpretability} to~\ref{sec:ablation_establishes}
bound the three analyses whose readings are most easily over-extended, the
external estimate, the treatment-independent label, and the feature ablation.
Section~\ref{sec:implications} draws out what the findings imply for how
studies of this kind should be reported, and
Section~\ref{sec:limitations} sets out the limitations in full.


\section{The Core Finding}
\label{sec:discussion_core}

The most important finding is definitional, not numerical: under a
conjunction-based onset rule the Sepsis-3 label on MIMIC-IV \emph{cannot}
contain an event before clinical action begins (Proposition~\ref{prop:h3}), and
under the standard $\min(t_{\text{susp}}, t_{\text{SOFA}})$ rule applied to the
identical stays it contains \pcPreTreatOnsetsVarF{} of them at
\pcPreTreatAurocPrimaryF{} AUROC. A published MIMIC-IV sepsis result is
therefore only as informative about physiological deterioration as its onset
rule permits, and a result reported without the timing branch is measuring
anticipation of clinician behaviour inside a window that behaviour opened.
The Utility Score stays close to the no-alert strategy for every model because
the temporal split exposes them to a test population where the same
treatment-related signals do not rescue performance as consistently as in the
training years, and because the score charges for every false-alert hour while
crediting only the first correct one. AUROC survives because the ordinal
ranking of patients is preserved even as the achievable alert volume becomes
unusable.

That divergence is what Wang et al.\ \cite{wang2025srpm} documented across
ninety-one models moving from internal to external validation, and what
Section~\ref{sec:auroc_illusion} identified as the field's central evaluation
failure. The result here locates it one step earlier. On MIMIC-IV, under an
honest time-based split, discrimination and clinical value have already parted
company \emph{internally}: there was scarcely any margin left for external
validation to destroy.


\section{H1 and H6 Together}
\label{sec:discussion_h1h6}

H1 (label dominance) and H6 (model-class null) are in productive tension, as
designed, and the resolution is sharper than either hypothesis anticipated,
because both resolve in the same direction, which the pre-registration treated
as unlikely.

H6 is \emph{not} rejected: the two flexible estimators are equivalent to within
\pcHSixEst{} AUROC, the whole interval inside the 0.02 bound. H1 is also not
confirmed, and by a wide margin in the direction opposite to the one
hypothesised: the model-class spread \pcSpreadModelClass{} exceeds the label
spread \pcSpreadLabel{} by \pcHOneEst, with an interval
[\pcHOneCiLo, \pcHOneCiHi] that excludes zero.

Read together these are not in conflict; they locate the model-class effect
precisely. The large model-class spread is the distance between \emph{families}
of estimator, a six-item bedside score against a fitted hazard model, not
between one fitted model and another. Within the family of flexible estimators
the choice is free (H6), and across families it is the single largest factor in
the decomposition (H1). The conventional literature, which fixes the label and
compares flexible models to each other, is therefore measuring the one contrast
that has almost no room left in it, while the choice that would move the number (fit anything flexible at all, rather than deploy a score) is not in
dispute and the choice that moves the cohort, the label, is invisible to
the metric being reported.

The H6 result also settles, for this task and this dataset, the narrower
version of the statistics-versus-machine-learning question left open in
Section~\ref{sec:false_dichotomy}. Christodoulou et al.'s
\cite{christodoulou2019cox} equivalence finding was established for binary
outcomes on tabular data, and Fagerstr{\"o}m et al.\ \cite{fagerstrom2019lisep}
showed it breaking on the temporal early-warning axis, where a deep model
detected patients up to twenty hours earlier than a Cox baseline. Here, with
the label, the split, the feature set and the evaluation held fixed, an
interpretable discrete-time hazard model and gradient-boosted trees are
statistically equivalent, and the interpretable model is marginally ahead on
point estimates. That is not a contradiction of Fagerstr{\"o}m et al.: their
comparison was against a static Cox specification, which is exactly the
comparator this study retains as the deliberately weak one, and which loses by
a wide margin here too (\pcAurocCoxB{} against \pcAurocPrimaryB). The
equivalence holds between the primary model and the boosted trees; it does not
hold between either of them and the field's default survival specification.

This is also the clearest illustration in the thesis of why the inference layer
was necessary. On point estimates, H1 confirmed and the study told a tidy story
in which the label dominates. With intervals attached, that story is not
supportable, and the defensible finding is the more interesting one: two
sources of variance of comparable size, only one of which the field routinely
reports.


\section{What the External Estimate Can Support}
\label{sec:external_interpretability}

Discrimination is estimated on \pcExtNEventsA{}, \pcExtNEventsB{} and
\pcExtNEventsC{} external events under Variants A, B and C
(Section~\ref{sec:external_validation}). The interpretability floor adopted
here is \pcExtMinEvents{} events, following the external-validation sample-size
guidance of Riley et al.\ \cite{riley2019sample} and Collins et al.\
\cite{collins2024tripodai}: below it the confidence interval on an external
AUROC is wider than the degradation the study is trying to detect.
\textbf{All three variants fall below it.} The degradations are accordingly
stated as directions, not as magnitudes, and the H5 interval is the formal
expression of that limit.

Two further quantities are unavailable at this event count in any case. AUPRC is
dominated by sampling noise at a few dozen positives, and the Utility Score
requires enough septic stays for the timeliness reward to be estimated at all.
Neither is interpreted externally.

\paragraph{The antibiotic-only arm.} Relaxing the culture requirement across
the full cohort raises the event count by two orders of magnitude (\pcExtAbxNEventsA, \pcExtAbxNEventsB{} and \pcExtAbxNEventsC{} events across
\pcExtAbxNRows{} person-hours), and therefore clears the floor with room to
spare. Under that anchor the primary model scores \pcExtAbxAurocPrimaryA,
\pcExtAbxAurocPrimaryB{} and \pcExtAbxAurocPrimaryC.

This arm labels a broader and less specific target: treated suspicion of
infection, without the organ-dysfunction conjunction's culture limb. An AUROC
computed against it therefore answers a different question and is never
compared like-for-like with the internal Sepsis-3 estimates, nor pooled with the
primary arm. What it does supply is a consistency check that the primary arm
cannot supply for itself. The degradation it shows is of the same sign, the same
label ordering and a somewhat larger magnitude than the Sepsis-3 arm's, on a
sample large enough to resolve it. Two estimates that disagree by two orders of
magnitude in event count and agree on direction are better evidence of a
direction than either is alone, which is the strongest statement the external
analysis in this thesis supports, and deliberately weaker than a claim about how
much discrimination is lost.

Against the published record assembled in Section~\ref{sec:generalisation},
the direction is unsurprising and the magnitude is unresolvable. Moor et al.\
\cite{moor2023review} report an external penalty of 0.085 AUROC across four
international databases; the point estimate here, \pcHFiveEst, is about
three-quarters of that. Where this study cannot follow Moor et al.\ is in
attributing the loss, because their harmonised databases shared labelling
inputs and eICU-CRD does not.


\section{What Variant D Does and Does Not Establish}
\label{sec:variant_d_establishes}

Variant D's high AUROC must be read with the limitation stated in advance in
Section~\ref{sec:alt_labels}. Its label is a threshold function of recorded
physiology, and five of the six SOFA components are themselves model inputs, so
its discrimination is partly self-fulfilling. The \emph{level} of
\pcAltAurocFullPrimaryD{} (\pcAltAurocFullGbtD{} for GBT) is not evidence that
deterioration prediction is a solved problem, and it is not comparable with
Variant B's \pcAurocPrimaryB.

Three things it does establish:

\begin{enumerate}
    \item \textbf{The label sensitivity is not wholly definitional.} At
    $\kappa = \pcKappaVsBD$ and a \pcPctExactMatchD\% onset match, a
    treatment-independent label
    disagrees with Sepsis-3 about both membership and timing. Had D closely
    tracked B, the H1 result would have been an artefact of how one definition
    is read; it is not.

    \item \textbf{The alert burden is label-invariant; the utility ceiling is
    not.} Both alternative labels are markedly more scoreable than Sepsis-3.
    The best achievable Utility Score rises from \pcUtilityMaxGbtB{} under
    Variant B to \pcUtilityMaxGbtE{} under E and \pcUtilityMaxGbtD{} under D,
    with the primary model following the same pattern
    (\pcUtilityMaxPrimaryE{} and \pcUtilityMaxPrimaryD). What does not improve
    to a deployable level is the alarm volume: even D's best configuration
    carries \pcUtilityOptBurdenGbtD{} false alerts per true alert, better than
    Variant B's \pcUtilityOptBurdenGbtB{} but still nine times the 1.4
    benchmark of Moor et al.\ \cite{moor2023review}. Contribution C3 survives
    the change of label in its alarm-volume form, which the pre-registered
    variants alone could not show, and does not survive in the form of a
    utility ceiling near zero. The rise is itself expected from the caveat
    above: a label built from recorded physiology is easier to alert on than
    one built from a clinician's decision.

    \item \textbf{The model-class ordering is label-invariant.} GBT leads on
    AUROC under all three labels and, once each model is evaluated at its own
    best threshold, on the Utility Score as well. A reversal appears only if
    every model is scored at the same 0.5 cut-off, which rewards a model for
    being too poorly calibrated to reach it. The
    metric-choice finding (H5) rests on the gap between discrimination and
    achievable utility, not on a re-ordering of models.
\end{enumerate}


\section{What the Ablation Establishes}
\label{sec:ablation_establishes}

The ablation is a small experiment and should not be asked to carry more than
it can. It is exploratory, it is run on one cohort, and with \pcHTwoNTerms{}
covariates the stability counts on either side are small enough that the
comparison is descriptive rather than tested.

What it does establish is a distinction the thesis needs and could not
otherwise make. The argument that the sepsis label is constituted by clinician
action invites an obvious objection: that the same is true of everything in an
ICU dataset, so the observation is unfalsifiable and therefore empty. The
ablation answers that objection empirically. The feature set divides into
predictors that record the patient and predictors that record the clinician,
the division is made in advance rather than fitted, and the two behave
differently under a change of label. That is a testable version of the claim,
and it is tested here rather than asserted.

It could have failed, and the form of the result matters as much as its
direction. Had discrimination fallen substantially, the finding would have been
that the models lean on clinician behaviour: supportive of the same argument,
but a concern about what reported performance is measuring, belonging as much in
the limitations as in the results. Had stability been unchanged, the
concentration of instability in the action-derived features would have been
small-sample noise and the reading would have had to be withdrawn. Neither
happened. Discrimination is essentially unaffected
(\pcAblDeltaAurocPrimaryB{} for the primary model under the primary label) while
the unstable count falls from \pcHTwoEst{} to \pcAblHTwoNUnstable, which is the
one outcome that separates the two hypotheses cleanly: features contributing
almost nothing to discrimination cannot be carrying the model's predictive
weight, and features whose removal takes five of six unstable terms with them
are where the label sensitivity lives.

The reading this licenses is bounded and worth stating exactly. It is
\emph{not} that the action-derived features are unimportant: a feature that
moves AUROC by less than 0.01 may still be the reason a particular alert fires.
It is that on this cohort the discrimination such a model reports does not
depend on them, while its cross-label stability does. A published effect
estimate for a treatment indicator from a model of this kind is therefore
reporting something about the labelling regime as much as about the patient, and
the cheapest available remedy, dropping the explicit action features, costs
almost nothing in the metric the field reports.

The corresponding limitation is stated in Section~\ref{sec:limitations}: the
retained laboratory values are themselves forward-filled from tests somebody
chose to order, so the boundary drawn here is conservative rather than clean. A
model with no trace of clinician behaviour in it is not constructible from this
data, and the ablation measures the effect of removing the explicit traces, not
all of them.


\section{Implications for How These Studies Are Reported}
\label{sec:implications}

The preceding sections bound what this study establishes. This one states what
follows for the design of the next one, because the findings bear less on which
model to build than on what a sepsis prediction study has to report before its
headline number can be read at all.

\paragraph{A label variant is not a sensitivity analysis.}
Section~\ref{sec:label_object} observed that the literature demonstrates
\emph{that} the label matters without attaching an interval to how much. The
result here sharpens the reason that gap is costly. Label choice moves AUROC by
\pcSpreadLabel{} and cohort membership by far more, the two narrowest
readings agree at $\kappa = \pcKappaVsBA$ at near-identical prevalence, so a
study that reports one definition's AUROC has reported a number that would
barely move under a different definition while every count, rate and
denominator beneath it moved substantially. The consequence is that reporting a
second variant as a robustness check answers the wrong question. What needs
reporting is the agreement between the variants, because that is where the
variation is, and it is invisible in the metric the field publishes.

\paragraph{An onset rule is a claim about what can be predicted, and belongs in
the methods.}
Section~\ref{sec:label_action} recorded Kamran et al.'s finding that models
partly learn to predict clinician behaviour, and this study locates the
mechanism one level down: not in the feature set but in the timestamp. Two
studies can share a cohort, a definition and a model and still differ on
whether pre-onset prediction is measurable, purely through the choice between
$t_{\text{susp}}$ and $\min(t_{\text{susp}}, t_{\text{SOFA}})$. That choice is
routinely left implicit in a sentence about Sepsis-3. It should be stated
explicitly, with its consequence for the evaluation window stated alongside it,
because under a conjunction rule an empty pre-treatment window is not evidence
of anything (Proposition~\ref{prop:h3}).

\paragraph{Discrimination and clinical value must both be reported, and neither
substitutes.}
Section~\ref{sec:auroc_illusion} set out why AUROC conceals the failure Wang et
al.\ documented: it is monotone-invariant and dominated by true negatives.
This study adds the observation that a single threshold conceals it a second
way. At the event rate here the conventional 0.5 cut-off separates models by
calibration rather than by clinical value, and the primary model's zero at that
cut-off is an artefact of never alerting there. A threshold-dependent metric is
interpretable only as a swept ceiling with the attaining operating point
reported beside it, and the alert burden at that point (\pcUtilityOptBurdenGbtB{} false alerts per true one here, against the 1.4
benchmark of Moor et al.\ \cite{moor2023review}) is the number that decides
deployability. It appears in almost none of the ninety-one models Wang et al.\
\cite{wang2025srpm} reviewed.

\paragraph{An external cohort should be characterised before it is scored.}
Section~\ref{sec:generalisation} assembled the published external penalties
into a single frame and treated them as comparable. The experience here is that
they are not automatically so. The binding constraint on the external estimate
in this study is not sample size but whether the label's inputs exist in the
destination database at all, and that constraint is invisible in a reported
AUROC and in a reported cohort size alike. A study transporting a
treatment-anchored label should report the coverage of each labelling input in
the destination cohort before it reports any performance number, and should
state the event count the estimate rests on rather than the stay count.

\paragraph{Model-class comparisons are the least informative axis available.}
Section~\ref{sec:false_dichotomy} argued that the statistics-versus-machine-learning
partition does not carve the field. The decomposition supports a stronger
version: within the family of flexible estimators the choice is free to within
\pcHSixEst{} AUROC, so a study comparing one flexible model against another is
measuring the axis with the least room left in it, while the label choice that
sets its cohort goes unreported. The reporting recommendation follows directly: fit any flexible estimator, report which one and stop defending the choice,
and spend the space on the label and the onset rule instead.

\paragraph{A subgroup analysis needs an interpretability floor declared in
advance.}
Section~\ref{sec:equity_lit} noted that the nearest prior work measures
subgroup disparity for mortality among already-septic patients rather than for
onset. Extending it to onset prediction exposed a problem that is structural
rather than particular to this cohort: at a low event rate most demographic
strata carry too few events for a subgroup AUROC to mean anything, and an
interval reaching the chance line is easily read as a null result rather than
as an absent measurement. Declaring the floor before estimation, tabulating the
below-floor strata rather than suppressing them, and excluding them from the
corrected family is what keeps the distinction visible. Without it, a study
underpowered for a fairness question can report one as answered.


\section{Limitations}
\label{sec:limitations}

\begin{enumerate}
    \item \textbf{Single-centre development.} All models were developed on
    MIMIC-IV, which originates from a single academic medical centre (Beth
    Israel Deaconess Medical Center, Boston). Clinical practices, patient
    demographics, and EHR documentation patterns at BIDMC may not represent
    those at other institutions. The equity and label-sensitivity findings in
    particular should not be assumed to transfer to hospitals with different
    patient populations or antibiotic prescribing cultures. The eICU-CRD
    analysis is the only external evidence available here, and the limitation
    below bounds how much it can carry.

    \item \textbf{External validation is bounded by label transport, not by
    cohort size.} eICU-CRD contains \pcExtNCohortStays{} unit stays, but the
    Sepsis-3 suspicion anchor requires a microbiology culture and
    \texttt{microLab} documents one for only \pcExtMicroCoverage\% of the
    analysis cohort. The primary external analysis is confined to that
    sub-cohort, because in the remaining stays no onset is observable under any
    variant; the resulting \pcExtNEventsB{} events under the primary label sit
    below the \pcExtMinEvents{}-event interpretability floor
    (Section~\ref{sec:external_interpretability}). Adding person-hours from the
    uncovered stays would enlarge the denominator without adding a single
    observable onset, so this is not a limitation additional data of the same
    kind would remove, it would require a destination database whose
    microbiology recording matches MIMIC-IV's. The consequence is specific and
    should not be over-read in either direction: the study establishes that
    discrimination degrades on transport, and cannot establish by how much.
    External Utility Score and AUPRC are not interpreted at all. The
    antibiotic-only sensitivity arm is well powered and agrees on direction,
    but it labels a different target and cannot stand in for the Sepsis-3
    comparison.

    \item \textbf{The external degradation is not decomposed into its causes.}
    An AUROC that falls on transport may reflect case-mix differences,
    documentation and measurement-frequency differences, genuine differences in
    how sepsis presents across two hundred hospitals, or residual imperfection
    in the transported label. This study measures the total and attributes none
    of it. The antibiotic-only arm bounds the label's contribution loosely (degradation persists when the culture limb is removed, so it is not solely
    an artefact of culture coverage), but a proper attribution would require
    a destination database with equivalent labelling inputs, which is precisely
    what was unavailable.

    \item \textbf{The NEWS2 comparator is not scored identically in the two
    cohorts.} NEWS2 charges three points for a patient who is not alert and
    none for a patient whose consciousness was not assessed. In MIMIC-IV the
    flag is derived from the recorded Glasgow Coma Scale with an unmeasured
    value treated as ``not alert'', so the small minority of person-hours
    carrying no GCS after forward-fill are charged those three points; in
    eICU-CRD, where the variable is absent from the periodic vitals table
    altogether, an unmeasured value is propagated as unmeasured and charged
    nothing. AUROC is a ranking statistic and the internal effect is confined
    to that minority of rows, but the two cohorts' NEWS2 \emph{thresholds} are
    not on a common scale, so the internal and external specificity and
    alert-burden figures for that comparator should be read within a cohort and
    not across the pair. Nothing in the decomposition, the hypothesis verdicts
    or the primary-model results depends on this: NEWS2 enters the study as a
    deliberately weak reference point, and the internal-versus-external
    contrast that H5 tests is computed on the primary model.

    \item \textbf{The feature set cannot be fully purged of clinician action.}
    The ablation in Section~\ref{sec:ablation} removes the
    \pcAblNDropped{} predictors that record clinician behaviour explicitly, vasopressor exposure and the per-test order indicators, but the
    laboratory values it retains are forward-filled from tests somebody chose
    to order, and even a vital sign is recorded on a schedule that reflects how
    closely a patient is being watched. The boundary drawn is therefore
    conservative rather than clean: what is measured is the effect of removing
    the explicit traces of clinical attention, not of removing all of them. A
    predictor set with no such trace is not constructible from routinely
    collected ICU data, which is itself part of the thesis's argument rather
    than only a limitation of it.

    \item \textbf{Deep survival model not implemented.} A deep survival
    comparator was planned but not built. Fagerstr{\"o}m et al.\
    \cite{fagerstrom2019lisep} report a recurrent architecture detecting
    septic-shock onset up to twenty hours earlier than a Cox baseline on the
    same features and target, so the absence of a deep temporal comparator is
    the one place where H6's equivalence result is genuinely incomplete: it
    establishes equivalence between the primary model and gradient-boosted
    trees, not between the primary model and every flexible estimator. This
    does not affect the structural findings about the label.

    \item \textbf{Limited feature set.} Clinical notes, procedure codes, and
    ventilator settings were not included. Richer features might improve
    discrimination but would not change the treatment-constitutionality finding.

    \item \textbf{Part of the label-variant spread is definitional, not
    empirical.} Variants A--C are three readings of one definition and all
    three fire on the same culture-plus-antibiotic anchor. Sepsis-3 makes onset
    a function of treatment timing, so moving the antibiotic--culture windows
    \emph{must} move the onsets; some of the H1 label spread is guaranteed by
    construction and would appear even if the physiology were identical across
    patients. Section~\ref{sec:label_anchor_results} separates the two by
    labelling each limb of the Sepsis-3 conjunction on its own, but the
    separation is post-hoc (Protocol Amendment 2) and rests on two variants,
    not on a systematic survey of alternative definitions.

    \item \textbf{The primary model has no analytic standard errors.} A
    cluster-robust sandwich estimator is not defined for the fitted multinomial
    object, so no per-coefficient variance is available
    (Section~\ref{sec:primary_model_spec}). Every interval reported for the
    primary model comes from the stay-level cluster bootstrap, which respects
    the same clustering but cannot supply a coefficient-level standard error.
    H2 is reported as a descriptive count for this reason, and any effect
    estimate from this model should be read as a point estimate without an
    attached uncertainty.

    \item \textbf{The specification requires a penalty to be identified at all.}
    The unpenalised likelihood is completely separated on this design under
    every label variant (Section~\ref{sec:primary_model_spec}), so the reported
    fits are ridge-penalised. The penalty is small and no result depends on its
    value, test AUROC moves by less than 0.02 across a 20-point path, but
    the coefficients are shrunk estimates, not maximum-likelihood ones, and the
    H2 stability counts are properties of the penalised solution. Reporting an
    unpenalised fit here would have been worse, not better: its coefficients
    have no unique maximiser, so cross-label comparison of them is undefined.

    \item \textbf{Feature construction and data cleaning are single-analyst
    work.} The covariate panel is assembled from raw chart and lab events by
    one pipeline, with no independent re-implementation to check it against.
    Two classes of defect are known to be possible in that setting and are
    guarded rather than eliminated: a declared predictor silently failing to
    reach a model because of a name mismatch, and implausible recorded values
    entering the fit. The pipeline now errors on any declared-but-absent
    feature rather than dropping it, and applies an explicit plausibility range
    to all nineteen physiological variables before forward-fill: MIMIC-IV
    chartevents contains values several orders of magnitude outside the
    physiological range, and such rows carry unbounded leverage. Neither guard
    is a substitute for independent replication of the feature pipeline.

    \item \textbf{The pre-registered onset rule is not the standard one.}
    Variants A--C place the onset at $t_{\text{susp}}$ and use the SOFA
    criterion only to gate membership, whereas Seymour et al.\
    \cite{seymour2016assessment} and the PhysioNet/CinC 2019 Challenge
    \cite{reyna2020utility} place it at
    $\min(t_{\text{susp}}, t_{\text{SOFA}})$. The pre-registered family
    therefore cannot address any question about onset \emph{timing} on its own,
    H3 among them: under its rule an empty pre-treatment window is a theorem
    rather than a result (Proposition~\ref{prop:h3}). Variant F supplies the
    measurement the pre-registered design cannot, but it is post-hoc (Protocol
    Amendment 4), it uses the full SOFA score including the vasopressor terms, so its onsets are free of the antibiotic-and-culture anchor but not of
    treatment in general, and its $t_{\text{SOFA}}$ is searched only inside
    the 72-hour modelled panel, so a rise occurring earlier in the admission
    cannot be detected. A design that pre-specified both timing rules across all
    three variants would be stronger than one that adds the second afterwards.

    \item \textbf{The treatment-independent label is not a validated construct.}
    Variant D is an operational deterioration label built for this comparison,
    not a consensus definition with an evidence base behind it, and its onsets
    have no external adjudication. Its label is also a threshold function of
    recorded physiology of which five of six components are model inputs, so
    its discrimination \emph{level} is not comparable with Variant B's, only
    its disagreement with Variant B about which patients deteriorate, and when,
    carries weight.

    \item \textbf{Specific to Sepsis-3 and to MIMIC-IV.} These findings apply
    to the Sepsis-3 definition on MIMIC-IV. SEP-1 and the CDC Adult Sepsis Event
    were not implemented; both are defined on antibiotic days, so neither would
    have tested treatment-independence, but each would give a different cohort
    and might behave differently in other respects. Other databases may behave
    differently again.

    \item \textbf{Operating points are not externally recalibrated.} The
    thresholds in Section~\ref{sec:clinical_utility} are derived on the test
    set at which they are evaluated, so the reported PPV and false-alarm rates
    are optimistic: a prospectively fixed threshold would perform no better.
    The same applies to the maximum Utility Scores in
    Table~\ref{tab:utility_optimum}, which take the best of sixty thresholds on
    the test set and are therefore upper bounds rather than estimates of what a
    deployed system would achieve; they are labelled exploratory for that
    reason. Selecting the threshold on the training years and transferring it
    would be the honest deployment estimate, and would be lower.
    The external operating points in Table~\ref{tab:external_utility} are
    likewise re-derived on eICU rather than transferred from MIMIC-IV, which
    isolates the discrimination loss but hides the additional loss a
    transferred threshold would incur.

    \item \textbf{Alarm episodes are defined by contiguity, not by
    acknowledgement.} False-alarm episodes are counted as contiguous runs of
    alert hours. A real system would apply suppression, snoozing and
    escalation logic, all of which would reduce the episode counts in
    Table~\ref{tab:patient_level}. The alert-hour rates are the conservative
    bound and the episode rates the optimistic one; the true burden lies
    between them.

    \item \textbf{The subgroup discrimination analysis is underpowered, and its
    strata are administrative.} Two distinct limits apply and neither
    substitutes for the other. On precision: only \pcNSubgroupsAboveFloor{} of
    the racial strata carry at least \pcMinSubgroupEvents{} sepsis onsets, the
    floor below which an AUROC is reported but not interpreted here
    (Section~\ref{sec:equity}), and no contrast in family E1 survives
    false-discovery correction. That is a statement about precision, not about
    fairness: the data are consistent both with no disparity and with
    disparities larger than any effect this thesis reports elsewhere. On
    construct validity: MIMIC-IV's \texttt{race} field is a single
    administrative code assigned at registration, so the strata compared are
    documentation categories rather than populations. Broad and granular codes
    drawn from the same population appear as separate levels (``White''
    excludes stays coded ``White -- Other European'', and ``Asian'' excludes
    those coded ``Asian -- Chinese''), so the reference level is a residual
    category and the contrasts are not between disjoint groups. ``Unknown'' and
    ``Unable to obtain'' record the absence of an answer. The
    label-sensitivity result (Table~\ref{tab:equity_label}) is the better
    powered of the two equity analyses because it is a stay-level proportion
    rather than a discrimination metric, but it inherits the same construct
    limitation, and its two significant strata are precisely the two
    missingness categories, which is why it is read in this thesis as a
    finding about documentation completeness rather than about race.

    \item \textbf{The novelty claims rest on a narrative, not systematic,
    search.} The search protocol in Section~\ref{sec:search_strategy} is
    documented and every included source was verified against a primary record,
    but no title-and-abstract screening pass over a full de-duplicated yield was
    performed and no second reviewer screened independently. The ``to the best
    of our knowledge'' claims for contributions C1 and C5 are therefore bounded
    by that coverage: a systematic review could surface a prior study combining
    multi-variant label sensitivity with subgroup equity analysis, and this
    search would not necessarily have found it. Both claims are claims about a
    \emph{combination} of analyses; each component separately has clear
    precedent, which Section~\ref{sec:search_strategy} identifies.

    \item \textbf{Five protocol amendments are post-hoc.} The inference layer
    (intervals and multiplicity control, Section~\ref{sec:inference}), the
    alternative label limbs D and E (Section~\ref{sec:alt_labels}), the model
    identification and data-plausibility changes
    (Section~\ref{sec:primary_model_spec}), the standard onset timing of
    Variant F (Section~\ref{sec:variant_f}) and the action-derived feature
    ablation (Section~\ref{sec:ablation}) were all specified after the
    pre-registered analyses had been run and inspected. None changes an
    estimand, threshold, direction or hypothesis, and H1--H6 and families E1/E2
    are computed from unchanged inputs; but a reader should treat the verdicts
    as a pre-specified hypothesis set analysed under a retrospectively specified
    error-control plan, and the amendment results as exploratory. Amendment 4 in
    particular enlarges family E3, which tightens rather than relaxes the
    correction applied within it.

    \item \textbf{Two hypotheses admit no test statistic.} H2 is a stability
    criterion rather than a contrast, and the fitted multinomial carries no
    variance estimate; H3 is degenerate because its outcome is absent from the
    evaluation window \emph{by construction} rather than by chance, so no test
    of it could ever have been informative. Both are reported descriptively
    while still consuming a slot in the family size, which makes the correction
    on the remaining four hypotheses stricter, not weaker. For H1, H4 and H6
    the correction determines nothing, each carries an adjusted $p$-value of
    1.000 and each fails on its point estimate. H5 is the exception: its
    adjusted value is \pcHFivePHolm, and it fails on the width of its interval
    rather than on its direction.
\end{enumerate}
